\documentclass[paper=a4, fontsize=11pt]{scrartcl} % A4 paper and 11pt font size

\newcommand{\assignment}{5}
\newcommand{\duedate}{March 15, 2023}
\usepackage[top=1in, bottom=1.5in, left=1in, right=1in]{geometry}
\usepackage{fancyhdr} % Required for custom headers
\usepackage{lastpage} % Required to determine the last page for the footer
\usepackage{extramarks} % Required for headers and footers
\usepackage[usenames,dvipsnames]{color} % Required for custom colors
\usepackage{graphicx} % Required to insert images
\usepackage{listings} % Required for insertion of code
\usepackage{courier} % Required for the courier font
\usepackage{amsmath}
\usepackage[super]{nth}
\usepackage{booktabs}
\usepackage[usenames,dvipsnames]{xcolor}
\usepackage{tcolorbox}
\usepackage{tabularx}
\usepackage{array}
\usepackage{colortbl}

%\usepackage[T1]{fontenc} % Use 8-bit encoding that has 256 glyphs
%\usepackage{fourier} % Use the Adobe Utopia font for the document - comment this line to return to the LaTeX default
\usepackage[english]{babel} % English language/hyphenation
\usepackage{amsmath,amsfonts,amsthm} % Math packages
\usepackage{graphicx}

\usepackage{hyperref}
\hypersetup{
  colorlinks   = true, %Colours links instead of ugly boxes
  urlcolor     = blue, %Colour for external hyperlinks
  linkcolor    = blue, %Colour of internal links
  citecolor   = red %Colour of citations
}

\usepackage{fancyhdr} % Custom headers and footers
\pagestyle{fancyplain} % Makes all pages in the document conform to the custom headers and footers
\fancyhead{} % No page header - if you want one, create it in the same way as the footers below
\fancyfoot[L]{} % Empty left footer
\fancyfoot[C]{} % Empty center footer
\fancyfoot[R]{\thepage} % Page numbering for right footer
\renewcommand{\headrulewidth}{0pt} % Remove header underlines
\renewcommand{\footrulewidth}{0pt} % Remove footer underlines
\setlength{\headheight}{13.6pt} % Customize the height of the header
\newcommand{\ts}{\textsuperscript}

\numberwithin{equation}{section} % Number equations within sections (i.e. 1.1, 1.2, 2.1, 2.2 instead of 1, 2, 3, 4)
\numberwithin{figure}{section} % Number figures within sections (i.e. 1.1, 1.2, 2.1, 2.2 instead of 1, 2, 3, 4)
\numberwithin{table}{section} % Number tables within sections (i.e. 1.1, 1.2, 2.1, 2.2 instead of 1, 2, 3, 4)

\setlength\parindent{0pt} % Removes all indentation from paragraphs - comment this line for an assignment with lots of text

% Default fixed font does not support bold face
\DeclareFixedFont{\ttb}{T1}{txtt}{bx}{n}{8} % for bold
\DeclareFixedFont{\ttm}{T1}{txtt}{m}{n}{8}  % for normal

%----------------------------------------------------------------------------------------
%	CODE BLOCKS
%----------------------------------------------------------------------------------------

\usepackage{adjustbox}
\usepackage{listings}
\usepackage{color}

\definecolor{dkgreen}{rgb}{0,0.6,0}
\definecolor{gray}{rgb}{0.5,0.5,0.5}
\definecolor{mauve}{rgb}{0.58,0,0.82}

\lstdefinelanguage{Dockerfile}
{
  morekeywords={FROM, RUN, CMD, LABEL, MAINTAINER, EXPOSE, ENV, ADD, COPY,
    ENTRYPOINT, VOLUME, USER, WORKDIR, ARG, ONBUILD, STOPSIGNAL, HEALTHCHECK,
    SHELL},
  morecomment=[l]{\#},
  morestring=[b]"
}

\lstset{
    columns=flexible,
    aboveskip=5mm,
    belowskip=5mm,
    keepspaces=true,
    showstringspaces=false,
    basicstyle=\ttfamily,
    commentstyle=\color{gray},
    keywordstyle=\color{purple},
    stringstyle=\color{green}
}



%----------------------------------------------------------------------------------------
%	TITLE SECTION
%----------------------------------------------------------------------------------------

\usepackage{eso-pic}
% \usepackage[demo]{graphicx}
\newcommand\AtPageUpperRight[1]{\AtPageUpperLeft{%
   \makebox[\paperwidth][r]{#1}}}

\newcommand{\horrule}[1]{\rule{\linewidth}{#1}} % Create horizontal rule command with 1 argument of height

\title{	
\normalfont \normalsize
\textsc{Northeastern University,  Khoury College of Computer Science} \\ [25pt] % Your university, school and/or department name(s)
\horrule{0.5pt} \\[0.4cm] % Thin top horizontal rule
\huge CS 6220  Data Mining \textemdash~Assignment \assignment \\ % The assignment title
\Large \textbf{Due: \duedate (100 points)} % The assignment title
\horrule{2pt} \\[0.5cm] % Thick bottom horizontal rule
}

% Original in the document
\AddToShipoutPictureBG*{%
  \AtPageUpperRight{\raisebox{-\height}{\includegraphics[width=3cm]{include/logo}}}}

\author{
    \textbf{YOUR NAME} \\ 
    \textbf{YOUR GIT USERNAME} \\ 
    \textbf{YOUR E-MAIL}
}% INFORMATION

\date{} % Today's date or a custom date

\author{
    \textbf{YOUR NAME} \\ 
    \textbf{YOUR GIT USERNAME} \\ 
    \textbf{YOUR E-MAIL}
}% INFORMATION

\begin{document}

\maketitle % Print the title

{\huge \textbf{Na\"ive Bayes, Bayes Rules}} \\

The original performance of \href{https://www.ncbi.nlm.nih.gov/pmc/articles/PMC6339026/}{acoustic classification for Parkinsons Disease} leverages speech recordings from controlled subject responses from variety of questions. The task in the competition was to detect whether or not a person $X$ had Parkinsons disease from a sampling of data. As of 2018, the state of the art classifiers have achieved 90\% correct classification on a held out dataset, both for subjects who had Parkinsons and those who did not (at equal rates). So, when classifier $Y$ sees person $X$, it works correctly 90\% of the time. \\

{\Large \textbf{Question 1} [25 pts total]} \\
\\
Let's say that we run a clinic. This clinic leverages this classifier, which has 90\% accuracy. Also, let us say that we know that our current patient load is that 10\% of the population have Parkinsons and 90\% of the population do not. Let's also say that we're seeing patient $X$, and the classification algorithm has detected that they have Parkinson's disease. What's the probability that indeed $X$ has Parkinson's disease? \\

Come up with the numerical solution, and show your written work.


\vspace{1.5cm}
{\huge \textbf{Gradient Descent - Logistic Regression}}
\vspace{3mm}

We'll explore logistic regression with some data from Kaggle today. If you've done the lab in class, then this simply becomes an exercise of data transformations and preprocessing. (Since there is only one label that you are predicting, you may also have to ensure that the dimensions are correct.) \\

Recall that the logistic regression with binary cross-entropy cost function appears in the following forms: \\

\begin{equation}
    \mathcal{L}(W, b) = -\sum_i y_i \log{ h_{W,b}(\textbf{x}) } + (1 - y_i) \log{( 1 - h_{W,b}(\textbf{x})) } \nonumber
\end{equation}
where 
\begin{equation}
h_{W,b}(\textbf{x}) = \sigma \left( W^T\textbf{x} + b \right) \nonumber
\end{equation}
and
\begin{equation}
\sigma(z) = \frac{1}{1 - \exp(-z)} \nonumber
\end{equation} \\
\\
Also, recall that you can solve the logistic regression problem entirely through the gradients as provided below:

\begin{itemize}
    \item The gradient of $\mathcal{L}$ with respect to $W$: $\nabla_W \mathcal{L} = \left(\textbf{y} - h_{W,b}(\textbf{x}) \right) \cdot \textbf{x}^T$
    \item The gradient of $\mathcal{L}$ with respect to $b$: $\nabla_b \mathcal{L} = \left(\textbf{y} - h_{W,b}(\textbf{x}) \right)$
\end{itemize}

\vspace{1cm}
{\Large \textbf{Question 2a.)}} \\

Prove that the derivative of the sigmoid function is $\nabla_z \sigma(z) = \sigma(z) \left(1 - \sigma(z)\right)$. 


\vspace{1cm}
{\Large \textbf{Question 2b.)}} \\

The sinking of the Titanic is one of the most infamous shipwrecks in history. On April 15, 1912, during her maiden voyage, the widely considered “unsinkable” RMS Titanic sank after colliding with an iceberg. Unfortunately, there were not enough lifeboats for everyone onboard, resulting in the death of 1502 out of 2224 passengers and crew. While there was some element of luck involved in surviving, it seems some groups of people were more likely to survive than others. In this challenge, Kaggle asks you to build a predictive model that answers the question: “what sorts of people were more likely to survive?” using passenger data (i.e., name, age, gender, socio-economic class, etc). \\

Some tips and tricks:

\begin{itemize}
    \item For features with multiple values, try making one-hot encoded data. For example, number of siblings, you can try to have a five column array. \\
    
    \begin{tabular}{|c|c|c|c|c|}
        \hline
         No Siblings & 1 Sibling & 2 Siblings & 3 Siblings & 4+ Siblings  \\
         \hline
                0    &    1      &      0     &       0  &    0\\
        \hline
    \end{tabular} \\
    \item Play around with your learning rate. If your learning is too high, then you won't converge to the right answer. I used X features with a learning rate of.
    \item You needn't use all the features. Just because they're there, doesn't mean they're always useful.
\end{itemize}

You should adopt the code from the lecture (and not use a toolbox). See if you can do better than I can with a logistic regression. On the test set, I reached 79.2\% accuracy at a threshold of 0.5. Plot out the loss curve and print your overall accuracy.

\vspace{1cm}
{\Large \textbf{Question 2c.)}} \\

Using a toolbox (like Keras), try doing the same with a few more layers (i.e., a neural network). Diagram your neural network architecture and print out your overall \textbf{test set} accuracy. Make sure you do \emph{not} train on any samples from your test set.


\end{document}
